\partstart{Open Correspondence Program}\label{part:correspondence}\subsection*{Outlook}
\label{sec:outlook}\label{sec:sonoluminescence}\label{sec:fieldlimit}\label{sec:cmb}
A field of transducers must first supply a consistent global update: operations sharing registers must compose without conflicting writes. Its collective \emph{eigenmodes} are eigenvectors or invariant subspaces of a declared global operator. They are distinct from the binary Mode defined in Section~\ref{sec:mode}.

Balanced release retains its definition as efficient conversion of internal excitation into admitted propagating eigenmodes with complete information accounting. Its earlier illustrative output-weight formula and the conditional thermal ensemble are preserved in Appendix~\ref{app:collective}. Large excitation and efficient outgoing coupling are different operator properties. Temperature additionally requires an energy assignment and justified occupation statistics.

Sonoluminescence remains a downstream boundary-flux comparison, where a native release criterion would have to explain measured behavior beyond established stability and coupling constraints \citep{hilgenfeldt1996}. A CMB application additionally requires electromagnetic eigenmodes, energy calibration, and thermal and cosmological history. Neither example supplies a temperature, a topology, or evidence for a microscopic carrier by analogy. The shared field-limit requirement is a sequence of admitted models with specified observables and controlled convergence.


The Maxwell translation in Section~\ref{sec:maxwelltranslation} supplies a conventional reference calculation, while a native electromagnetic reduction remains a concrete target: derive admitted coupled channels with consistent divergence constraints, curl evolution, charge continuity and calibrated propagating solutions, rather than assign the WXYZTI letters to electric and magnetic components \citep{feynmanmaxwell}. Appendix~\ref{app:horizon} also preserves the $V_4^2\cong\F_2^4$ reference-fiber and $3+1$ proposals as conditional construction questions. Neither the integer $960$ nor a choice of one binary direction supplies a spacetime metric or replaces the certified time register.


The intermediate test is a source-selected admitted update that preserves its declared invariant and constraint surface on a common state domain. Sources must be represented as exchanges with specified complementary variables; a rotor analogue must identify an axis orbit and its stabilizer. Only after these finite construction gates should the controlled rotational or electromagnetic limit be tested. The periodic domain used for the Maxwell calculation is an analytic reference choice, not evidence for the global topology of a physical universe.
